\documentclass[10pt,twocolumn]{article}

% ── Page geometry ─────────────────────────────────────────────────────────────
\usepackage[a4paper, top=2.5cm, bottom=2.5cm, left=1.8cm, right=1.8cm,
            columnsep=0.8cm]{geometry}

% ── Core packages ─────────────────────────────────────────────────────────────
\usepackage{amsmath,amssymb}
\usepackage{graphicx}
\usepackage{booktabs}
\usepackage{array}
\usepackage{tabularx}
\usepackage{enumitem}
\usepackage{hyperref}
\usepackage{microtype}
\usepackage{abstract}
\usepackage{titlesec}
\usepackage{fancyhdr}
\usepackage{xcolor}
\usepackage{mdframed}
\usepackage{listings}
\usepackage{courier}
\usepackage{caption}
\usepackage{multicol}
\usepackage{parskip}

% ── Colors ────────────────────────────────────────────────────────────────────
\definecolor{navyblue}{RGB}{31,56,100}
\definecolor{steelblue}{RGB}{46,117,182}
\definecolor{lightblue}{RGB}{214,228,247}
\definecolor{codebg}{RGB}{245,245,245}
\definecolor{darkgray}{RGB}{89,89,89}
\definecolor{alertbg}{RGB}{255,248,220}
\definecolor{alertborder}{RGB}{204,153,0}

% ── Hyperref ──────────────────────────────────────────────────────────────────
\hypersetup{
  colorlinks=true,
  linkcolor=steelblue,
  urlcolor=steelblue,
  citecolor=steelblue
}

% ── Section heading styles ────────────────────────────────────────────────────
\titleformat{\section}
  {\large\bfseries\color{navyblue}}
  {\thesection.}{0.5em}{}
  [\vspace{-0.3em}\textcolor{steelblue}{\rule{\linewidth}{0.6pt}}]

\titleformat{\subsection}
  {\normalsize\bfseries\color{steelblue}}
  {\thesubsection}{0.5em}{}

\titleformat{\subsubsection}
  {\normalsize\bfseries\color{darkgray}}
  {\thesubsubsection}{0.5em}{}

\titlespacing*{\section}{0pt}{1.4ex plus .2ex}{0.8ex plus .1ex}
\titlespacing*{\subsection}{0pt}{1.0ex plus .1ex}{0.5ex plus .1ex}
\titlespacing*{\subsubsection}{0pt}{0.8ex}{0.3ex}

% ── Abstract style ────────────────────────────────────────────────────────────
\renewcommand{\abstractnamefont}{\normalfont\bfseries\color{navyblue}}
\renewcommand{\abstracttextfont}{\normalfont\small\itshape}
\setlength{\absleftindent}{0pt}
\setlength{\absrightindent}{0pt}

% ── Header / Footer ───────────────────────────────────────────────────────────
\pagestyle{fancy}
\fancyhf{}
\fancyhead[L]{\small\textcolor{darkgray}{Network Security --- Certificate Authority Chain}}
\fancyhead[R]{\small\textcolor{darkgray}{\thepage}}
\renewcommand{\headrulewidth}{0.4pt}
\renewcommand{\footrulewidth}{0.4pt}

% ── Listings (code blocks) ────────────────────────────────────────────────────
\lstset{
  basicstyle=\footnotesize\ttfamily,
  backgroundcolor=\color{codebg},
  frame=single,
  rulecolor=\color{steelblue!40},
  breaklines=true,
  breakatwhitespace=true,
  numbers=none,
  xleftmargin=4pt,
  xrightmargin=4pt,
  aboveskip=4pt,
  belowskip=4pt,
}

% ── Custom environments ───────────────────────────────────────────────────────
\newmdenv[
  linecolor=steelblue,
  linewidth=1pt,
  backgroundcolor=lightblue!40,
  innerleftmargin=6pt,
  innerrightmargin=6pt,
  innertopmargin=5pt,
  innerbottommargin=5pt,
  skipabove=6pt,
  skipbelow=6pt,
]{infobox}

\newmdenv[
  linecolor=alertborder,
  linewidth=1.2pt,
  backgroundcolor=alertbg,
  innerleftmargin=6pt,
  innerrightmargin=6pt,
  innertopmargin=5pt,
  innerbottommargin=5pt,
  skipabove=6pt,
  skipbelow=6pt,
]{alertbox}

% ── List spacing ──────────────────────────────────────────────────────────────
\setlist[itemize]{noitemsep, topsep=2pt, leftmargin=*}
\setlist[enumerate]{noitemsep, topsep=2pt, leftmargin=*}

% ── Paragraph spacing ─────────────────────────────────────────────────────────
\setlength{\parindent}{0pt}
\setlength{\parskip}{4pt}

% ─────────────────────────────────────────────────────────────────────────────
\title{%
  \vspace{-1.5em}
  {\LARGE\bfseries\color{navyblue}
   Certificate Authority Chain and Deployment\\[4pt]
   of a CA on a Linux Server}\\[6pt]
  {\normalsize\color{darkgray}
   Network Security Project --- Academic Report}
}
\author{%
  \textbf{Fernandez Ojeda Franklin 000541971} \\
  \textbf{Lecky-Thompson William 000613648} \\
  \textbf{Otto Aleksandra 000569128} \\
  \textbf{Rafaat Moheeb Eskandar Abanoub 000567614} \\
  \textbf{Vanderslagmolen Moïra 000547486} \\
}
\date{\small\today}

% ─────────────────────────────────────────────────────────────────────────────
\begin{document}

\twocolumn[
  \maketitle
  \vspace{-0.5em}
  \begin{onecolabstract}
    \noindent
    This report presents the design, deployment, and validation of a
    complete Public Key Infrastructure (PKI) on a Linux server using
    OpenSSL. Starting from the theoretical foundations of X.509 certificates
    and PKI trust models, it covers the full lifecycle of certificate
    management: creation of a two-tier Certificate Authority hierarchy
    (Root CA and Intermediate CA), issuance of server and client
    end-entity certificates, certificate revocation via CRL, and the
    deployment of Mutual TLS (mTLS) for bidirectional authentication.
    The practical implementation demonstrates how mTLS strengthens
    security beyond standard TLS by requiring both parties to present
    valid certificates during the handshake, and examines its relevance
    in modern Zero Trust architectures and microservice environments.
    \smallskip
  \end{onecolabstract}
  \vspace{1em}
  \hrule
  \vspace{1em}
]

% ─────────────────────────────────────────────────────────────────────────────
\section{Introduction}

\subsection{Objectives}

This report aims to achieve the following:

\begin{itemize}
  \item Explain the structure and role of X.509 digital certificates in PKI.
  \item Describe the components and trust models of Public Key Infrastructure.
  \item Present the certificate signing and verification process in detail.
  \item Provide a practical guide to deploying a root CA and intermediate CA
        on a Linux server using OpenSSL.
  \item Demonstrate how to issue end-entity certificates for TLS services.
  \item Implement and validate Mutual TLS (mTLS) for bidirectional
        authentication between client and server.
\end{itemize}

\subsection{Motivation}

Secure communication on the Internet depends fundamentally on a trustworthy
framework for managing digital identities. HTTPS, VPNs, email signing, and
code signing all rely on X.509 certificates and Certificate Authorities. Yet
the inner workings of PKI remain opaque to many practitioners. Beyond
standard server authentication, the growing adoption of Zero Trust
architectures makes client-side certificate authentication, as provided
by mTLS, increasingly critical in modern deployments.

\subsection{Report Structure}

\begin{enumerate}
  \item \textbf{Background:} X.509 certificates and the PKI framework.
  \item \textbf{Certificate Signing \& Verification:} full lifecycle and
        verification checks.
  \item \textbf{CA Deployment on Linux:} setting up root and intermediate CAs,
        issuing end-entity certificates, and configuring revocation.
  \item \textbf{Mutual TLS:} theory, handshake mechanics, and use cases.
  \item \textbf{Limitations of PKI:} structural and operational weaknesses
        illustrated by documented real-world incidents.
\end{enumerate}

% ─────────────────────────────────────────────────────────────────────────────
\section{Background}

\subsection{X.509 Certificates}

An X.509 certificate is a digital document that binds a public key to a
verified identity (a person, server, organization, or device). The
certificate is signed by a Certification Authority (CA), allowing relying
parties to trust its contents without prior contact with the owner.

The X.509 standard defines three versions: v1, v2, and v3. Version~3 is
universally adopted because it supports extensions that convey critical
security metadata. Table~\ref{tab:cert-fields} summarizes the core fields.

\begin{table}[h]
\centering
\caption{Core X.509 v3 Certificate Fields}
\label{tab:cert-fields}
\renewcommand{\arraystretch}{1.2}
\begin{tabularx}{\columnwidth}{>{\bfseries\footnotesize}p{2.1cm}
                                >{\footnotesize}X}
\toprule
Field & Description \\
\midrule
Version        & Format version (v1/v2/v3). \\
Serial Number  & Unique identifier assigned by the CA. \\
Sig.\ Algorithm& Algorithm used to sign (e.g.\ SHA-256/RSA). \\
Issuer Name    & Distinguished Name (DN) of the issuing CA. \\
Validity Period& The window of time in which the certificate is valid and trusted. \\
Subject Name   & DN of the certificate owner. \\
Public Key Info& Owner's public key and algorithm. \\
Extensions     & Key Usage, SAN, Basic Constraints, etc. \\
Signature      & CA's cryptographic signature over TBSCert (the certificate body). \\
\bottomrule
\end{tabularx}
\end{table}

\subsubsection{Version 3 Extensions}

Version 3 extensions are one of the main reasons X.509 certificates are useful
in real systems. The basic certificate fields identify the subject, issuer,
and public key, but the extensions tell software \emph{how} the certificate
may be used and what limits apply to it. Each extension has an identifier, a
value, and a \emph{critical} flag. If an application does not understand a
critical extension, it must reject the certificate rather than guess its
meaning \cite{rfc5280}.

In practice, extensions let a CA describe whether a certificate is for a web
server, a client, an intermediate CA, or a root CA. They also help software
build chains correctly, check revocation information, and apply naming or
policy restrictions. For this reason, extensions are not just optional extra
data; they are central to certificate validation \cite{rfc5280}.

Some of the most important Version~3 extensions are the following:

\begin{itemize}
  \item \textbf{Key Usage:} This extension lists the cryptographic operations
        allowed for the key, such as \texttt{digitalSignature},
        \texttt{keyEncipherment}, \texttt{keyCertSign}, and
        \texttt{cRLSign}. For example, a CA certificate should allow
        certificate signing, while a normal TLS server certificate should not.
  \item \textbf{Extended Key Usage (EKU):} EKU gives a more precise purpose to
        the certificate, such as \texttt{serverAuth}, \texttt{clientAuth},
        \texttt{codeSigning}, or \texttt{emailProtection}. In practice, both
        Key Usage and EKU should be consistent with the role of the
        certificate.
  \item \textbf{Subject Key Identifier (SKI) and Authority Key Identifier
        (AKI):} These two extensions help software link certificates together.
        SKI identifies the subject's key, while AKI points to the key of the
        issuing CA. They are especially useful when building and checking a
        certificate chain.
  \item \textbf{Certificate Policies:} This extension lists policy object
        identifiers (OIDs) that describe the rules under which the certificate
        was issued. Related extensions such as Policy Constraints and Inhibit
        Any Policy help control policy processing in more complex PKI
        environments.
  \item \textbf{Subject Alternative Name (SAN):} SAN can contain DNS names, IP
        addresses, email addresses, URIs, and other identities. In modern TLS,
        SAN is the main place checked for hostname matching. RFC~5280 also
        notes that if the subject field is empty, the SAN extension must be
        present and marked critical \cite{rfc5280}.
  \item \textbf{Basic Constraints:} This extension states whether a
        certificate belongs to a CA by using \texttt{CA:TRUE} or
        \texttt{CA:FALSE}. It may also include
        \texttt{pathLenConstraint}, which limits how many subordinate CA
        levels may appear below that certificate.
  \item \textbf{Name Constraints:} This extension is mainly used in CA
        certificates to restrict the names that subordinate certificates may
        contain. For example, a CA can be limited to issuing certificates only
        for a specific domain or namespace.
  \item \textbf{CRL Distribution Points and Authority Information Access
        (AIA):} These extensions tell clients where to find revocation and
        issuer information. CRL Distribution Points give the location of the
        certificate revocation list, while AIA may point to issuer
        certificates or OCSP services used during validation.
\end{itemize}

Overall, Version~3 extensions make certificate validation more precise and
more secure. They allow software to enforce the intended role of a
certificate, distinguish between CA and end-entity certificates, and obtain
supporting information such as revocation status and issuer details. Without
these extensions, modern PKI would be much harder to manage safely
\cite{rfc5280}.

Certificates are typically written in one of two formats:
\begin{itemize}
    \item \textbf{Distinguished Encoding Rules (DER):} a binary encoding format.
    \item \textbf{Privacy Enhanced Mail (PEM):} Base64-encoded DER with header
          and footer lines, commonly used in configuration files.
\end{itemize}

\subsection{Public Key Infrastructure (PKI)}

PKI is the overarching framework for managing the issuance, distribution,
validation, and revocation of digital certificates. Beyond cryptographic
algorithms, PKI defines the policies, roles, and procedures that allow parties
to establish trust across untrusted networks.

\subsubsection{PKI Trust Models}

\begin{itemize}
  \item \textbf{Hierarchical (Single Root):} a single root CA delegates trust
        downward to subordinate CAs. Used by public web PKIs (e.g.\ browser
        root stores).
  \item \textbf{Cross-Certification (Mesh):} multiple independent CAs establish
        trust relationships by mutually issuing certificates to one another.
        Instead of relying on a single root authority, trust is distributed
        among participating organizations.
  \item \textbf{Bridge CA:} a Bridge Certification Authority acts as a central
        hub that connects multiple independent PKI domains. Rather than every
        CA cross-certifying with every other CA, each participating PKI only
        establishes trust with the bridge CA.
\end{itemize}

\subsubsection{Main PKI Components}

A Public Key Infrastructure (PKI) consists of several components that work
together to establish, distribute, validate, and revoke digital trust.

\begin{itemize}
  \item \textbf{Certificate Authority (CA):} the trusted entity responsible for
        issuing, digitally signing, renewing, and revoking certificates. The
        CA binds a subject's identity to its public key, thereby acting as the
        root of trust within the PKI hierarchy.

  \item \textbf{Registration Authority (RA):} an entity that performs identity
        verification on behalf of the CA. The RA validates information provided
        by certificate applicants before forwarding approved requests to the
        CA for certificate issuance.

  \item \textbf{Digital Certificates:} electronic documents conforming
        typically to the X.509 standard. A certificate contains identifying
        information about the subject, the subject's public key, validity
        period, certificate serial number, issuer information, and the CA's
        digital signature.

  \item \textbf{Certificate Signing Request (CSR):} a structured request
        generated by an applicant when requesting a certificate. The CSR
        contains the applicant's public key and identifying information
        (such as domain name or organization details) and is digitally signed
        using the applicant's private key to prove possession of that key.

  \item \textbf{Certificate Repository:} a centralized directory service
        (commonly LDAP-based) used to store and distribute issued certificates,
        CA certificates, and Certificate Revocation Lists (CRLs). It allows
        relying parties to retrieve certificates required for validation.

  \item \textbf{Revocation Services:} mechanisms used to determine whether a
        certificate is still trusted before its expiration date. Common methods
        include Certificate Revocation Lists (CRLs) and the Online Certificate
        Status Protocol (OCSP).

  \item \textbf{Certificate Policy (CP) / Certification Practice Statement
        (CPS):} formal documents defining the PKI's security requirements,
        trust model, operational procedures, identity verification processes,
        and certificate management practices.
\end{itemize}

\subsubsection{PKI Operation Sequence}

The certificate issuance and validation lifecycle within a PKI typically
follows the sequence below:

\begin{enumerate}
  \item Generate an asymmetric cryptographic key pair consisting of a public
        key and a private key.
  \item Create a Certificate Signing Request (CSR) containing the public key
        and subject identity information. The CSR is digitally signed using the
        private key to prove ownership and possession of the corresponding key.
  \item Submit the CSR to a Registration Authority (RA) or directly to the CA.
  \item The RA verifies the applicant's identity according to the PKI policy
        and forwards the validated request to the CA.
  \item The CA validates the request and digitally signs the certificate using
        its private key, thereby binding the applicant's identity to the public
        key.
  \item The signed certificate is issued and delivered to the requester, who
        installs it on a server, device, application, or user system.
  \item The certificate may also be published to a certificate repository so
        that relying parties can retrieve it when necessary.
  \item During communication, relying parties validate the certificate by
        checking the CA signature, certificate chain, validity period,
        hostname or subject information, and revocation status using CRL or
        OCSP services.
  \item If a certificate is compromised, expired, or no longer trusted, the CA
        revokes it and updates revocation information accordingly.
\end{enumerate}

% ─────────────────────────────────────────────────────────────────────────────
\section{Certificate Signing and Verification}

The X.509 PKI framework, as standardized in ITU-T Recommendation X.509 and
RFC~5280 \cite{rfc5280}, establishes a hierarchical trust model wherein
digital certificate signing serves as the fundamental mechanism for
authenticating cryptographic identities. Through the application of asymmetric
cryptography, Certification Authorities generate unforgeable digital signatures
that cryptographically bind entity identities to their respective public keys.
This process, which involves the generation of a signed ASN.1-encoded
TBSCertificate structure, provides both data origin authentication and message
integrity through mathematical verification of the CA's signature using
standardized algorithms such as RSA-PKCS\#1~v1.5 or ECDSA \cite{rfc5280}.

The verification process constitutes a critical component of secure
communication protocols, including the TLS handshake \cite{rfc8446}. During
certificate chain validation, relying parties must not only verify the digital
signature's mathematical correctness but also enforce temporal validity
constraints, key usage policies, and path validation rules as specified in
RFC~5280 Section~6. This multi-faceted verification process ensures that only
properly authorized entities can participate in secure communications, while
effectively mitigating well-known attacks such as man-in-the-middle (MITM)
interposition.

\subsection{Certificate Signing Process}

\subsubsection{Key Pair Generation}

Before creating a CSR, the end-entity must generate an asymmetric key pair
using cryptographic software. The private key is kept secret to sign documents
and decrypt data, while the public key is included in the CSR. Two principal
algorithms are used:

\begin{itemize}
  \item \textbf{RSA:} used for both encryption and digital signatures.
        Recommended key sizes are 2048\,bits (minimum) or 4096\,bits for
        longer-lived certificates.
  \item \textbf{ECDSA:} a signature-only algorithm offering equivalent security
        with shorter keys. Curve P-256 is standard for TLS certificates
        \cite{rfc5280}.
\end{itemize}

Some good practices can be adopted to ensure the secure storage of the private
key:

\begin{itemize}
  \item Restrictive file permissions on Linux.
  \item Key encryption with a strong passphrase.
\end{itemize}

\subsubsection{CSR Structure (PKCS\#10)}

A Certificate Signing Request is the formal process of applying for a digital
certificate from a Certificate Authority. The CSR structure is defined in
PKCS\#10 \cite{pkcs10}, a security standard produced by RSA Data Security
Inc. It consists of a subject (DN info), the subject public key and its
algorithm, and a set of attributes. These fields are encoded using ASN.1
(Abstract Syntax Notation~1) and DER (Distinguished Encoding Rules), which
are binary encoding standards used in X.509.

A CSR contains three main parts:

\begin{itemize}
  \item \textbf{Certification Request Information:} contains pertinent
        information such as the public key, Common Name, Organization Name,
        Organizational Unit, Country, key type, and key size.
  \item \textbf{Signature Algorithm Identifier:} specifies the algorithm used
        to sign the CSR.
  \item \textbf{Digital Signature:} proves possession of the private key. The
        requester signs the CSR with their private key, allowing the CA to
        verify they control it.
\end{itemize}

The information provided in the CSR must be accurate and match the existing
domain registration. Otherwise, certificate issuance delays or even rejections
may occur.

\subsubsection{Submission and Identity Validation}

Once the CSR is completed, it is submitted to a trusted Certificate Authority,
which can be either internal (e.g.\ a corporate PKI) or external (e.g.\
Let's Encrypt or DigiCert). The CA performs identity verification, and the
exact process depends on the type of certificate requested:

\begin{itemize}
  \item \textbf{Domain Validation (DV):} the CA verifies only that the
        applicant controls the domain name. Verification is usually automated
        through an email challenge, a DNS TXT record, or an HTTP file upload.
        DV certificates are fast and inexpensive to issue, but they do not
        verify the identity of the organization behind the domain.

  \item \textbf{Organization Validation (OV):} the CA verifies both domain
        ownership and the legal existence of the organization. This may include
        checking business registries, verifying the organization's name,
        address, and registration in official databases, and conducting phone
        verification with authorized personnel. OV certificates provide
        stronger identity assurance than DV certificates.

  \item \textbf{Extended Validation (EV):} the CA performs strict identity
        verification, including legal, physical, and operational checks of the
        organization. There must be confirmation that the person requesting the
        certificate is authorized to act on behalf of the organization. EV
        certificates provide the highest traditional level of trust and are
        commonly used by banks, governments, and large organizations.
\end{itemize}

\subsubsection{Certificate Issuance}

Once validation passes, the CA constructs the TBSCertificate structure in
ASN.1, populates all fields and extensions, computes a cryptographic hash, and
signs it with the CA private key using the chosen algorithm
(e.g.\ SHA-256/RSA or ECDSA P-256). The complete CSR structure is digitally
signed by the CA using the CA's private key and a chosen signature algorithm.
The signed certificate is then returned to the requester in PEM or DER format.

\subsection{Certificate Verification Process}

To ensure that a digital certificate is valid, the client must proceed through
a multi-step verification process \cite{rfc5280}. This involves signature
validation, chain verification, and revocation status checks. Skipping any
step opens the door to attacks such as MITM or certificate substitution.

\subsubsection{Basic Certificate Checks}

\begin{itemize}
  \item \textbf{Validity Period:} a certificate contains a \texttt{Not Before}
        and a \texttt{Not After} timestamp. The verifier must check that the
        certificate is within its active date range.
  \item \textbf{Subject and Issuer Name:} the Subject field must match the
        entity, and the Issuer field must match the expected Certificate
        Authority.
  \item \textbf{Key Usage and Extended Key Usage (EKU):} the verifier must
        ensure that the certificate is used for its intended purpose. For
        example, a TLS server certificate should have \texttt{Server
        Authentication} in its EKU.
\end{itemize}

\subsubsection{Digital Signature Verification}

Every certificate is digitally signed by its issuing CA. The signature proves
that the CA has validated the certificate's contents. First, the certificate's
signed data and signature are extracted. The issuer's public key is then used
to verify the signature mathematically. If the signature is valid, it confirms
that the certificate was issued by that CA and that its content was not
modified since it was signed.

\subsubsection{Certificate Chain Validation}

Most certificates are not directly signed by a root CA but by intermediate
CAs. The verifier must validate the entire chain of trust \cite{rfc5280}.

The process starts with the end-entity certificate belonging to the
website or server. The verifier retrieves the issuer's certificate, or
intermediate CA certificate, and repeats this process until a trusted root CA
is reached. For each certificate in the chain, the verifier must:

\begin{enumerate}
  \item Verify the digital signature.
  \item Check the validity period.
  \item Ensure no revocation (via CRL or OCSP).
  \item Confirm basic constraints (e.g.\ intermediate CA must have
        \texttt{CA:TRUE}).
  \item Confirm the root CA exists in the client's trust store.
\end{enumerate}

\begin{alertbox}
\footnotesize\textbf{Security note:} A chain is only as strong as its weakest
link. If an intermediate CA is compromised, all certificates below it must be
revoked. This is why root CAs are kept \emph{offline} and intermediate CAs
handle day-to-day operations.
\end{alertbox}

\subsubsection{Revocation Status Checking}

Even if a certificate is properly signed, it may have been revoked due to
compromise or expiration. Two main revocation methods exist:

\begin{itemize}
  \item \textbf{CRL (Certificate Revocation List):} a list of revoked
        certificates published by the CA. The verifier downloads this CRL and
        checks if the certificate's serial number is listed.
  \item \textbf{OCSP (Online Certificate Status Protocol):} a real-time query
        to the CA to check certificate status \cite{rfc6960}. It returns
        \texttt{Good}, \texttt{Revoked}, or \texttt{Unknown}. If a certificate
        is found in the CRL or marked as revoked, it is not trusted.
        \emph{OCSP Stapling} allows the server to cache and present a signed
        OCSP response, eliminating client round-trips and improving privacy.
\end{itemize}

\subsubsection{Hostname Verification}

For HTTPS and other TLS connections, the certificate's Common Name (CN) or
Subject Alternative Name (SAN) must match the domain being accessed
\cite{rfc2818}. For example, if a user visits \texttt{https://example.com}
and the certificate is issued for \texttt{example.org}, the connection will
be rejected. Modern clients rely primarily on the SAN extension for this
check, as the legacy CN field is no longer used for hostname matching in
current browsers.

\subsection{Role in TLS Secure Communication}

X.509 certificates are used to secure communications between a client (such
as a web browser) and a web server. When a user connects to a website using
HTTPS, the web server presents a digital certificate compliant with the X.509
standard. This certificate includes information about the server's identity
(domain name), its public key, and a digital signature issued by a Certificate
Authority. The client then verifies the validity of the certificate by walking
up the certificate chain to a pre-installed trusted root authority in the
operating system or browser.

During a TLS~1.3 handshake \cite{rfc8446}, the server transmits its
certificate chain immediately after the Server Hello. The client performs all
the checks described above, and only on full success does it proceed to derive
session keys. This mechanism provides:

\begin{itemize}
  \item \textbf{Authentication:} the user can be sure they are communicating
        with the legitimate website and not an impostor.
  \item \textbf{Confidentiality:} data exchanged cannot be read or intercepted
        by third parties.
  \item \textbf{Integrity:} any alteration in the transmitted messages can be
        detected.
\end{itemize}

As such, X.509 certificates are essential for securing web traffic, preventing
attacks such as man-in-the-middle and DNS spoofing, and establishing trust
between users and online services.

% ─────────────────────────────────────────────────────────────────────────────
\section{CA Deployment on a Linux Server}

\subsection{Architecture Overview}

A standard deployment uses a \textbf{two-tier hierarchy}:

\begin{itemize}
  \item \textbf{Root CA} --- kept offline on an air-gapped machine or HSM.
        Only used to sign the intermediate CA certificate.
  \item \textbf{Intermediate CA} --- online, used for day-to-day certificate
        issuance. Compromise of the intermediate CA does not expose the
        root key.
\end{itemize}

\begin{infobox}
\footnotesize
\textbf{Best practice:} generate and store the root CA private key on an
air-gapped machine or HSM. Bring it online \emph{only} to sign intermediate
CA certificates. Store backups in separate physical locations.
\end{infobox}

\subsection{Directory Structure}

A well-organized directory structure is essential for auditability. The
working directory for the root CA should contain the following:

\begin{itemize}
  \item \texttt{certs/} --- issued certificates.
  \item \texttt{crl/} --- Certificate Revocation Lists.
  \item \texttt{newcerts/} --- archive of all signed certificates.
  \item \texttt{private/} --- private keys, with restricted permissions.
  \item \texttt{index.txt} --- the certificate database.
  \item \texttt{serial} --- the next serial number to be assigned.
  \item \texttt{openssl.cnf} --- the root CA configuration file.
  \item \texttt{intermediate/} --- the intermediate CA, with the same layout.
\end{itemize}

\subsection{Step-by-Step Deployment}

\subsubsection{Step 1 --- Create the Root CA}

The first step is to generate the root CA private key, protected with
AES-256 encryption. A self-signed root certificate is then created with a
long validity period (typically 20 years), using the \texttt{v3\_ca}
extensions defined in the configuration file. The resulting certificate
should be verified to confirm that all fields and extensions are correct.

\subsubsection{Step 2 --- Create the Intermediate CA}

An intermediate CA key is generated in the same way as the root key. A CSR
is then created for the intermediate CA and submitted to the root CA for
signing. The root CA signs the intermediate certificate using the
\texttt{v3\_intermediate\_ca} extensions, which include
\texttt{CA:TRUE} and an appropriate \texttt{pathLenConstraint}. Finally,
a certificate chain bundle is assembled by concatenating the intermediate
certificate and the root certificate into a single file, which is provided
to relying parties during TLS handshakes.

\subsubsection{Step 3 --- Issue End-Entity Certificates}

A server key and CSR are generated for the target service. The CSR must
include the Subject Alternative Name (SAN) extension specifying the domain
names or IP addresses the certificate will cover. The intermediate CA signs
the certificate using the \texttt{server\_cert} extensions, which set
\texttt{CA:FALSE}, restrict key usage to \texttt{digitalSignature} and
\texttt{keyEncipherment}, and set the Extended Key Usage to
\texttt{serverAuth}. The signed certificate, private key, and chain bundle
are then deployed to the web server (e.g.\ Nginx or Apache).

\subsubsection{Step 4 --- Configure Revocation (CRL \& OCSP)}

An initial CRL is generated and published at the URL specified in the
\texttt{crlDistributionPoints} extension of issued certificates. When a
certificate must be revoked, it is marked as revoked in the CA database and
the CRL is regenerated and republished. For real-time revocation queries, an
OCSP responder is deployed and its URL is included in the
Authority Information Access (AIA) extension of issued certificates.

\subsection{Key OpenSSL Configuration}

The \texttt{openssl.cnf} file controls all CA behavior. The
\texttt{[ v3\_ca ]} section defines the extensions applied to CA
certificates, including \texttt{CA:TRUE}, key usage for certificate and CRL
signing, and key identifier extensions. The \texttt{[ server\_cert ]}
section defines end-entity certificate extensions, setting \texttt{CA:FALSE},
restricting key usage, and requiring a \texttt{subjectAltName} entry.

\begin{alertbox}
\footnotesize\textbf{SAN is mandatory.} Modern browsers reject certificates
that lack a Subject Alternative Name, even if the Common Name matches.
Always include \texttt{subjectAltName} in server certificate configurations.
\end{alertbox}

% ─────────────────────────────────────────────────────────────────────────────
% MUTUAL TLS
% ─────────────────────────────────────────────────────────────────────────────
\section{Mutual TLS (mTLS)}

\subsection{Definition and Motivation}

Standard TLS provides \emph{one-way} authentication: the server presents a
certificate to the client, which verifies it against a trusted root CA. The
client, however, remains anonymous to the server. This is sufficient for
public-facing websites, where the server does not need to know \emph{who}
is connecting---only that the communication is confidential.

Mutual TLS (mTLS) extends this model to \emph{two-way} authentication.
Both the server and the client present X.509 certificates, and both verify
each other's certificate against a trusted CA before the connection is
established. Neither party can impersonate the other, and an attacker who
intercepts the channel cannot participate without a valid certificate issued
by the trusted CA.

This bidirectional trust is particularly relevant in environments where the
identity of the connecting entity, not just the confidentiality of the
channel, must be cryptographically guaranteed.

\subsection{mTLS Handshake Mechanics}

In a standard TLS~1.3 handshake \cite{rfc8446}, the server sends its
certificate chain after the \texttt{ServerHello} message. The client
verifies it and, on success, derives session keys. In mTLS, the server
additionally sends a \texttt{CertificateRequest} message, which signals
that the client must also authenticate. The extended handshake proceeds
as follows:

\begin{enumerate}
  \item The client sends \texttt{ClientHello}.
  \item The server responds with \texttt{ServerHello}, its certificate chain,
        and a \texttt{CertificateRequest}.
  \item The client verifies the server certificate against its trust store.
  \item The client sends its own certificate chain in response to the
        \texttt{CertificateRequest}.
  \item The client sends a \texttt{CertificateVerify} message, containing a
        signature over the handshake transcript produced with its private key.
        This proves that the client possesses the private key corresponding
        to the public key in its certificate.
  \item The server verifies the client certificate against its configured CA
        and checks the \texttt{CertificateVerify} signature.
  \item Both parties derive session keys and the connection proceeds.
\end{enumerate}

If either certificate fails validation, due to an invalid signature,
expired validity, revoked status, or missing \texttt{clientAuth} in the EKU
of the client certificate, the handshake is aborted and the connection is
rejected.

\begin{infobox}
\footnotesize
\textbf{Key extension:} client certificates used in mTLS must include
\texttt{clientAuth} in their Extended Key Usage field. In our deployment,
this is set in the \texttt{[ client\_cert ]} block of
\texttt{openssl.cnf}.
\end{infobox}

\subsection{Comparison: TLS vs.\ mTLS}

Table~\ref{tab:tls-mtls} summarizes the key differences between standard
TLS and mTLS.

\begin{table}[h]
\centering
\caption{TLS vs.\ mTLS comparison}
\label{tab:tls-mtls}
\renewcommand{\arraystretch}{1.2}
\begin{tabularx}{\columnwidth}{>{\bfseries\footnotesize}p{2.2cm}
                                >{\footnotesize}X
                                >{\footnotesize}X}
\toprule
 & \textbf{TLS} & \textbf{mTLS} \\
\midrule
Server auth   & Yes & Yes \\
Client auth   & No  & Yes (certificate) \\
Client anon.  & Yes & No \\
Handshake     & Standard & Extended (+CertReq) \\
Typical use   & Public websites & APIs, microservices, VPNs \\
Revocation    & Server cert & Both certs \\
\bottomrule
\end{tabularx}
\end{table}

\subsection{Relevance in Modern Network Security}

mTLS has become a foundational building block in several current security
architectures:

\begin{itemize}
  \item \textbf{Zero Trust Architecture:} the Zero Trust model assumes that
        no entity---inside or outside the network perimeter---is trusted by
        default. Every connection must be authenticated and authorized
        explicitly. mTLS provides the cryptographic identity layer that
        enables this, ensuring that even an attacker who has gained access
        to the internal network cannot establish connections without a valid
        certificate.

  \item \textbf{Microservice communication:} in distributed systems, services
        communicate over internal networks that may be shared or untrusted.
        Service meshes such as Istio and Linkerd use mTLS transparently
        between every pair of services, so that each service can verify it
        is talking to a legitimate peer and not a compromised component.

  \item \textbf{API security:} enterprise APIs increasingly require client
        certificates instead of---or in addition to---API tokens. Certificate-
        based authentication is stronger than bearer tokens because the private
        key never leaves the client and cannot be stolen from a database.

  \item \textbf{IoT device authentication:} in IoT deployments, each device
        is provisioned with a unique certificate at manufacture time. When
        the device connects to a backend server, mTLS ensures that only
        legitimate, enrolled devices are accepted, preventing rogue device
        injection.
\end{itemize}

\subsection{Certificate Validity and Renewal}

One operational challenge of mTLS is certificate lifecycle management.
Certificate authorities and browser requirements have progressively shortened
maximum validity periods: from several years before 2015, to 398 days by
2020, and proposals from major vendors now target 47 days or fewer
\cite{rfc5280}. Short-lived certificates limit the window of exposure if a
private key is compromised, but they require automated renewal processes to
avoid service disruption.

In our deployment, renewal is handled by a Bash script executed daily via
\texttt{cron}, the Linux task scheduler. The script is idempotent: it
inspects every certificate in its inventory and triggers a renewal only
when the remaining validity falls below a 30-day threshold, mirroring the
behaviour of Let's Encrypt's \texttt{certbot}. When renewal is required,
the script generates a fresh key pair and a new CSR, submits it to the
Intermediate CA, and verifies the resulting certificate against the chain
\emph{before} swapping any live files. The previous certificate and key
are first copied to a timestamped backup directory, and Nginx is reloaded
only after \texttt{nginx -t} confirms that the new configuration is valid.
If any of these checks fails, the previous certificate remains in place
and the service is not disrupted.

A separate weekly job regenerates the Certificate Revocation List well
before its 30-day \texttt{nextUpdate} boundary, since an expired CRL would
cause Nginx to reject every connection regardless of the validity of the
end-entity certificates. The two jobs are kept on independent schedules
because they concern different state---the end-entity certificate
inventory on one side, the revocation database on the other---and
decoupling them keeps each script focused on a single responsibility.

For unattended execution, the Intermediate CA passphrase is stored in a
root-owned file with \texttt{0600} permissions and supplied to OpenSSL via
the \texttt{-passin file:} option. This is a deliberate trade-off: a
passphrase file on the same machine reduces protection against host
compromise, but it allows the entire renewal pipeline to run without human
intervention. In a production environment, the Intermediate CA key would
be stored in a Hardware Security Module or a Key Management Service so
that the key material never leaves the secure boundary.

\begin{alertbox}
\footnotesize\textbf{Operational note:} in mTLS deployments, \emph{both}
the server and client certificates must be renewed before expiry. An
expired client certificate causes the same handshake failure as an invalid
one. The script therefore iterates over the complete certificate inventory
rather than only the server side, and each entry declares its own
extension block (\texttt{server\_cert} or \texttt{client\_cert}) so that
the correct Extended Key Usage is preserved across renewals.
\end{alertbox}

\subsection{Limitations of mTLS}

Despite its security benefits, mTLS introduces operational complexity that
must be carefully managed:

\begin{itemize}
  \item \textbf{Certificate distribution:} every client must be provisioned
        with a valid certificate before it can connect. In large deployments
        this requires a scalable enrollment process, often handled by an
        automated PKI management system.
  \item \textbf{Revocation propagation:} if a client certificate is
        compromised, it must be revoked and the CRL or OCSP responder
        updated promptly. Until the revocation is published, the compromised
        certificate remains usable.
  \item \textbf{Compatibility:} standard web browsers do not present client
        certificates automatically. mTLS is therefore suited to machine-to-
        machine communication rather than general user-facing web traffic.
  \item \textbf{Overhead:} the additional certificate exchange and verification
        steps add a small latency cost to connection establishment, though
        this is negligible for most applications.
\end{itemize}

% ─────────────────────────────────────────────────────────────────────────────
% LIMITATIONS OF THE PKI MODEL (NEW SECTION)
% ─────────────────────────────────────────────────────────────────────────────
\section{Limitations of the PKI Model}

Despite the strong cryptographic foundations of X.509 and the maturity of
modern PKI tooling, the global Certificate Authority ecosystem has suffered
recurring failures over the past fifteen years. The limitations discussed
below are not theoretical: each is illustrated by a well-documented
real-world incident, and each motivated specific design choices in our
own deployment. They show that the weak points of PKI lie not in its
mathematics but at its edges---operations, human processes, governance,
and revocation infrastructure.

\subsection{Technical Compromise of a CA}

A legitimate CA can be breached through poor operational security: outdated
servers, insufficient network segmentation, or weak monitoring. If a single
CA present in the browser trust store is compromised, an attacker can forge
\emph{valid certificates for any domain} (Google, banks, governments), and
victim browsers will accept them silently. This enables large-scale
man-in-the-middle attacks on HTTPS.

\textbf{Case study --- DigiNotar (2011).} The Dutch CA DigiNotar was
breached through an obsolete public-facing web server and insufficient
internal network segmentation. The attacker issued \textbf{531 fraudulent
certificates}, including a wildcard \texttt{*.google.com}, which was used
to intercept the Gmail traffic of approximately \textbf{300{,}000 Iranian
citizens} in what was likely a state-sponsored MITM operation. All major
browsers distrusted DigiNotar within weeks, and the company \textbf{filed
for bankruptcy in September 2011}\footnote{\url{https://en.wikipedia.org/wiki/DigiNotar}; see also Fox-IT, \textit{Operation Black Tulip: Report of the investigation into the DigiNotar Certificate Authority breach}, 2012.}.

This incident demonstrates that cryptographic strength offers no protection
against poor operational hygiene. It is precisely the reason our deployment
keeps the Root CA \emph{offline} and uses an online Intermediate CA only
for day-to-day operations.

\subsection{Commercial Pressure and ``Too Big to Fail'' CAs}

Public CAs are commercial enterprises that sell certificates at scale, and
this economic pressure can lead to lax validation, poorly supervised
resellers, and the dismissal of internal errors. When a CA becomes dominant
on the market, sanctioning it becomes politically and technically expensive:
a significant fraction of the web could be rendered inaccessible overnight.

\textbf{Case study --- Symantec distrust (2017--2018).} Symantec, together
with its VeriSign, Thawte, and GeoTrust brands, issued \textbf{more than
30\% of all valid web certificates} at the time. A Google Chrome
investigation initially identified 127 misissued certificates, but the
total grew to at least \textbf{30{,}000} over several years, many of them
emitted through unsupervised resellers. Google announced a progressive
distrust of Symantec certificates in Chrome 66 and Chrome 70, and Symantec
ultimately \textbf{sold its entire CA division to DigiCert for approximately
\$1 billion}\footnote{\url{https://www.thesslstore.com/blog/google-chrome-distrust-symantec-ssl-certificates/}; original announcement on the Chromium blog: \textit{Intent to Deprecate and Remove: Trust in existing Symantec-issued Certificates}.}.

This case illustrates that the commercial PKI model has a structural
conflict of interest, and that market concentration weakens the ecosystem's
ability to sanction misbehaviour.

\subsection{Revocation Is Technically Broken}

The two main mechanisms for checking whether a certificate has been revoked,
CRL and OCSP, have flaws that make them effectively unusable on the public
web:

\begin{itemize}
  \item \textbf{CRL} files list every revoked certificate and can grow to
        several megabytes, making them impractical to download on every
        connection.
  \item \textbf{OCSP} performs a real-time query to the CA, which is slow
        and leaks the user's browsing history to the CA \cite{rfc6960}.
  \item If the OCSP request fails---due to an outage or to an attacker
        actively blocking it---browsers \textbf{fail soft}: they accept the
        certificate anyway to avoid breaking the web. A network-positioned
        attacker can therefore neutralise revocation simply by dropping the
        OCSP request.
\end{itemize}

\textbf{Case study --- Heartbleed (April 2014).} The Heartbleed vulnerability
in OpenSSL exposed TLS private keys on approximately \textbf{17\% of HTTPS
servers worldwide}, forcing thousands of operators to revoke their
certificates. However, Chrome had already disabled OCSP by default, and its
alternative mechanism (CRLSets) contained only about \textbf{24{,}000
entries out of more than 2 million revoked certificates}---roughly
\textbf{98\% ignored}\footnote{\url{https://www.netcraft.com/blog/certificate-revocation-why-browsers-remain-affected-by-heartbleed/}; see also A.~Langley, \textit{Revocation doesn't work}, ImperialViolet blog, March 2011: \url{https://www.imperialviolet.org/2011/03/18/revocation.html}}. An
attacker who had stolen a private key could keep using it even after formal
revocation.

This connects directly to Section~4 (CRL configuration) and motivates the
broader industry trend of shortening certificate validity periods (398 days,
then 90 days, with 47 days now proposed), as well as the automatic renewal
script described in Section~5.

\subsection{Default Trust and CA Governance}

Modern browser trust stores contain 150--200 CAs that users trust
implicitly, without ever being asked to choose. Several of these CAs operate
under jurisdictions that can compel them to cooperate. Unlike the previous
case of an accidentally compromised CA, here the CA itself is the threat:
it may \emph{deliberately} issue fraudulent certificates, or be forced to,
and the bare X.509 model provides no native mechanism to restrict which CA
may issue for which domain.

\textbf{Case study --- CNNIC / MCS Holdings (2015).} CNNIC, China's official
CA, was present in every major trust store. CNNIC issued an
\textbf{intermediate certificate to MCS Holdings}, an Egyptian company with
no legitimate CA infrastructure, and \textbf{without any name constraint},
meaning MCS could issue certificates for arbitrary domains. MCS produced
\textbf{fake Google certificates} and installed them in a corporate firewall
to intercept its employees' encrypted traffic. The incident was detected
through \textbf{Certificate Transparency}, after which Google and Mozilla
distrusted CNNIC\footnote{\url{https://security.googleblog.com/2015/03/maintaining-digital-certificate-security.html}; see also \url{https://en.wikipedia.org/wiki/CNNIC\#Issuance_of_fake_certificates}}.

This case shows that ``default'' trust is a structural weakness of the
model. It directly motivated \textbf{Certificate Transparency} (RFC~6962),
which forces CAs to publish every issued certificate to public,
append-only logs, making large-scale misissuance detectable.

\subsection{Human Factor: Employees and Resellers}

CA security ultimately depends on humans: administrators, employees,
resellers. All cryptographic safeguards are bypassed if an administrative
account is phished or if a partner reuses credentials. In practice, attacks
on CAs almost always exploit the human factor rather than break the
underlying cryptography.

\textbf{Case study --- ComodoHacker (2011).} A solo attacker calling
themselves ``ComodoHacker'' compromised a \textbf{Comodo reseller account
based in Italy} by stealing valid credentials. There was no sophisticated
exploit and no cryptographic break---only a poorly protected account at a
business partner. The attacker issued \textbf{nine fraudulent certificates}
for Google, Yahoo, Skype, Mozilla, and login.live.com, and claimed the
operation was politically motivated, drawing comparisons with the later
DigiNotar attack\footnote{\url{https://en.wikipedia.org/wiki/Comodo_Cybersecurity\#Certificate_hacking}; see also Mozilla Security Blog, \textit{Comodo Certificate Issue Follow-up}, March 2011.}.

The modern response to this class of attack relies on mandatory multi-factor
authentication, strict reseller audits, and real-time monitoring of
certificate issuance, complemented by Certificate Transparency logs as a
public check.

\subsection{Operational Errors on the Operator Side}

Even in the absence of any attack, organisations routinely compromise
themselves: forgotten renewals, leaked private keys, and TLS
misconfigurations are common. The larger the deployment---hundreds or
thousands of certificates---the higher the probability of error.

\textbf{Case study --- Equifax (2017).} The US credit-rating agency Equifax
used an HTTPS-aware intrusion detection system that inspected encrypted
traffic for attacks. The internal certificate enabling this inspection
\textbf{expired and was not renewed for ten months}. During that period,
the detection tool was effectively blind, and attackers exfiltrated the
personal data of \textbf{145 million individuals} (Social Security numbers,
dates of birth, addresses) without triggering any alert. No CA was
breached, no cryptography was broken---a single forgotten certificate was
enough\footnote{US House Committee on Oversight and Government Reform, \textit{The Equifax Data Breach}, December 2018, p.~3: \url{https://republicans-oversight.house.gov/wp-content/uploads/2018/12/Equifax-Report.pdf}}.

This is precisely the failure mode that the \textbf{automatic renewal
script} described in Section~5 is designed to prevent, and it is also the
philosophy behind tools such as \texttt{certbot} for Let's Encrypt.

\begin{alertbox}
\footnotesize\textbf{Synthesis:} the recurring pattern across these
incidents is that PKI breaks at its edges---operational hygiene, human
processes, governance, and revocation infrastructure---rather than in its
mathematical core. The architectural choices in our deployment (offline
Root CA, automated renewal, scheduled CRL regeneration, restrictive file
permissions, strict EKU separation between server and client certificates)
are direct responses to these documented failure modes.
\end{alertbox}


% ─────────────────────────────────────────────────────────────────────────────
% AI STRATEGY SECTION
% ─────────────────────────────────────────────────────────────────────────────
\section{AI Strategy}

\subsection{Tools Used}

Throughout this project, the primary AI tool used was \textbf{Claude}
(Anthropic). It was used in three distinct roles: assistance with
implementation, support for report writing, and guidance for information
retrieval with verifiable sources.

\subsection{Implementation Assistance}

For the practical deployment of the CA infrastructure, mTLS configuration,
and the certificate renewal script, Claude was used as a technical assistant
rather than a code generator. The approach followed was to provide the AI
with the full context of the task---the existing directory structure, the
relevant \texttt{openssl.cnf} sections, and the precise objective---and to
ask for targeted help on specific steps, such as the correct flags for a
given OpenSSL command or the structure of a particular Nginx directive.

Every command and script produced with AI assistance was subsequently
executed and tested on the actual deployment environment. Understanding the
output was a requirement: if a command worked but its behaviour was not
fully understood, it was not accepted. This ensured that the implementation
reflects genuine understanding rather than blind execution of generated code.

\subsection{Report Writing}

Claude was also used to support the writing of this report. Its role was
limited to reformulation and grammar correction: sections were drafted by
the group, then submitted to the AI with instructions to improve clarity,
fix grammatical errors, or rephrase awkward sentences---while preserving
the original meaning and technical content. The ideas, structure, and
technical judgements in this report are entirely our own. Re-reading each
revised paragraph against the original draft ensured that no content was
altered beyond stylistic improvement and that no hallucinated or inaccurate
statements were introduced.

\subsection{Information Retrieval and Source Verification}

When researching topics such as the mTLS handshake mechanics, the evolution
of certificate validity periods, or the role of mTLS in Zero Trust
architectures, Claude was asked to provide not only explanations but also
references to primary sources: RFCs, official documentation, and peer-reviewed
material. Information obtained this way was cross-checked against the cited
sources before being included in the report. For example, the description
of the TLS~1.3 handshake in Section~5 was verified against RFC~8446
\cite{rfc8446}, and the X.509 extension semantics were checked against
RFC~5280 \cite{rfc5280}. Any claim that could not be traced to a verifiable
source was either removed or reformulated with appropriate hedging.

\subsection{Prompt Optimisation}

A key lesson learned early in the project was that the quality of AI output
is directly proportional to the precision and completeness of the prompt.
Vague or open-ended prompts tend to produce generic answers that require
significant correction. The approach adopted was therefore to make every
prompt as specific as possible, providing all relevant context upfront:
the current state of the deployment, the exact file or configuration
section in question, the error message or behaviour observed, and the
precise expected outcome. Rather than asking the AI to \emph{find} a
solution, prompts were structured so that the AI had all the information
needed to \emph{confirm or refine} a solution already partially formed.
This reduced the risk of hallucination and kept the AI's role firmly in the
position of assistant rather than decision-maker.

\begin{infobox}
\footnotesize
\textbf{Example --- imprecise prompt:}
\textit{``How do I configure mTLS?''}\\[4pt]
\textbf{Example --- precise prompt:}
\textit{``We have an Nginx server with a two-tier PKI (Root CA and
Intermediate CA built with OpenSSL). The client certificate was issued
using the \texttt{client\_cert} extension block with
\texttt{clientAuth} in EKU. What is the minimal Nginx configuration
to enforce mandatory client certificate verification using
\texttt{ssl\_verify\_client on}, pointing to our
\texttt{ca-chain.cert.pem} file, with a verify depth of 2?''}
\end{infobox}

\subsection{Contribution to Exam Preparation}

Beyond the practical tasks, Claude served as a personal tutor during the
study phase. Concepts encountered in the course slides that were not
immediately clear---such as the role of \texttt{pathLenConstraint},
the difference between CRL and OCSP Stapling, or the mechanics of the
mTLS handshake---were explored interactively. Unlike a static textbook,
the AI could adjust its level of explanation, provide concrete analogies,
and answer follow-up questions immediately. This made it a useful complement
to the lecture material when preparing for the individual questions expected
during the oral defence.

\subsection{Limitations and Critical Assessment}

The use of AI assistance also highlighted several limitations that required
active management:

\begin{itemize}
  \item \textbf{Hallucination risk:} the AI occasionally produced plausible
        but incorrect technical details, particularly for specific OpenSSL
        command flags or Nginx directive names. Systematic testing of every
        command in the actual environment was the primary mitigation.
  \item \textbf{Outdated information:} AI models have a knowledge cutoff and
        may not reflect the most recent changes in standards or tooling.
        For time-sensitive facts---such as current maximum certificate
        validity periods---primary sources and recent official documentation
        were always consulted.
  
\end{itemize}


% ─────────────────────────────────────────────────────────────────────────────
\begin{thebibliography}{9}

\bibitem{rfc5280}
D.~Cooper et al.,
\textit{Internet X.509 PKI Certificate and CRL Profile},
RFC~5280, IETF, May~2008.

\bibitem{rfc2818}
E.~Rescorla,
\textit{HTTP Over TLS},
RFC~2818, IETF, May~2000.

\bibitem{rfc6960}
S.~Santesson et al.,
\textit{Online Certificate Status Protocol --- OCSP},
RFC~6960, IETF, June~2013.

\bibitem{rfc8446}
E.~Rescorla,
\textit{The Transport Layer Security (TLS) Protocol Version 1.3},
RFC~8446, IETF, August~2018.

\bibitem{pkcs10}
RSA Laboratories,
\textit{PKCS \#10: Certification Request Syntax Standard},
Version~1.7, 2000.

\bibitem{itu509}
ITU-T,
\textit{Recommendation X.509: The Directory --- Public-Key and
Attribute Certificate Frameworks}, 2019.

\end{thebibliography}

\end{document}